Repeated rejection makes the final choice uniform among all noncollinear triples. Translate the fountain to the origin, so the lamps are the nonzero points of $\{-2,-1,0,1,2\}^2$.

First count collinear triples. In the full $5$ by $5$ grid, horizontal and vertical lines contribute $100$, slopes $\pm1$ contribute $40$, and slopes $\pm2,\pm\dfrac{1}{2}$ contribute $12$. No other slope supports three grid points. Removing the center removes respectively $12,12,4$ triples. Hence there are $124$ collinear triples and
\[
\binom{24}{3}-124=1900
\]
possible final triangles.

For the numerator, the $16$ rays from the origin alternate between eight rays with two lamps and eight with one. Opposite rays have the same size. A ray with $w$ lamps has $12-w$ lamps strictly within the next counterclockwise open half-plane. Every triple contained in an open half-plane has a unique first occupied ray in its angular span. Counting by that ray gives
\[
8\left[2\binom{10}{2}+\binom22\cdot10\right]
+8\binom{11}{2}=1240.
\]
Triples with the origin on an edge, but not all collinear, are counted by selecting an opposite-ray pair and a third point off its line. The four lines with two lamps on each ray give $4\cdot4\cdot20=320$; the four with one lamp on each ray give $4\cdot1\cdot22=88$. In addition, $4\binom43=16$ collinear triples have the origin in their segment hull.

The remaining triples strictly contain the origin, so their count is
\[
2024-1240-320-88-16=360.
\]
The probability is $\dfrac{360}{1900}=\dfrac{18}{95}$, giving $\boxed{m+n=113}$.
\begin{asy}
texpreamble("\usepackage{amsmath,amssymb}");
size(240);
defaultpen(fontsize(9pt));
for(int i=-2;i<=2;++i)for(int j=-2;j<=2;++j)if(i!=0||j!=0)dot((i,j));dot((0,0),red+6bp);draw((-2,-1)--(2,-1)--(0,2)--cycle,blue);label("fountain",(0,0),NE);
\end{asy}
